\documentclass[11pt,a4paper]{article}

\usepackage{fontspec}
\usepackage[a4paper,margin=2.0cm]{geometry}
\usepackage{graphicx}
\usepackage{tabularx}
\usepackage{array}
\usepackage{booktabs}
\usepackage{titlesec}
\usepackage{enumitem}
\usepackage{setspace}
\usepackage{caption}
\usepackage[table]{xcolor}
\usepackage[hidelinks]{hyperref}

\setmainfont{TeX Gyre Heros}
\defaultfontfeatures{Ligatures=TeX}
\fontsize{11}{14}\selectfont

\setlength{\parindent}{0pt}
\setlength{\parskip}{2pt}
\setlength{\tabcolsep}{4pt}
\renewcommand{\arraystretch}{1.08}
\raggedbottom
\setlength{\textfloatsep}{8pt plus 2pt minus 2pt}
\setlength{\floatsep}{8pt plus 2pt minus 2pt}
\setlength{\intextsep}{8pt plus 2pt minus 2pt}
\captionsetup{font=small,labelfont=bf,skip=2pt}
\titlespacing*{\section}{0pt}{5pt}{2pt}
\titlespacing*{\subsection}{0pt}{4pt}{2pt}
\setlist[itemize]{leftmargin=*,nosep}
\setlist[enumerate]{leftmargin=*,nosep}
\definecolor{coverblue}{HTML}{1F3A56}
\definecolor{coverline}{HTML}{C6D3DF}

\begin{document}

\begin{titlepage}
\thispagestyle{empty}
\centering
\vspace*{1.8cm}
{\Large\color{coverblue}\textbf{CS2SD}}\\[0.25em]
{\large Software Systems Design}

\vspace{1.2cm}
{\Huge\textbf{Stage 1 Report}}\\[0.45em]
{\Large Amazon Enterprise Online POS System}\\[0.2em]
{\large Use Case and Activity Design}

\vspace{0.9cm}
\color{coverline}\rule{0.82\textwidth}{1.2pt}
\color{black}

\vspace{0.9cm}
\renewcommand{\arraystretch}{1.16}
\setlength{\tabcolsep}{7pt}
\begin{table}[h]
\centering
\normalsize
\begin{tabularx}{\textwidth}{|>{\raggedright\arraybackslash}p{5.2cm}|X|}
\hline
\textbf{Module Code:} & CS2SD \\ \hline
\textbf{Assignment Report Title:} & Stage 1 for design of the software system: user requirements \\ \hline
\textbf{Student Number (e.g.\ 25098635):} & 32804061, 32803755, 32803720 \\ \hline
\textbf{Actual hrs spent for the assignment:} & Approx.\ 20 hours \\ \hline
\textbf{Which AI tools used (if applicable) and brief justification:} & AI-assisted wording refinement, report structuring, and consistency checking between the text and UML models; all scope decisions, diagram logic, and evaluation were produced and reviewed by the group. \\ \hline
\end{tabularx}
\end{table}
\vfill
{\small This report presents the Stage 1 OO analysis and design of an enterprise online POS system using UML use case and activity modelling.}
\end{titlepage}

\setcounter{page}{1}

\section{Introduction}

This report examines an Amazon enterprise online POS system designed to support guest browsing, registered purchasing, seller-side marketplace activity, and administrative supervision within a single business setting. The problem space is not only technical; it also involves social and organisational concerns, which makes the system a suitable subject for OOAD (Sommerville, 2016). A range of human and machine actors interact around the same transaction process, and design choices are shaped by issues such as privacy, accessibility, payment protection, and stock reliability (Booch, et al., 1999). Against that background, the report pursues two aims. The first is to identify actors, functional requirements, and operating constraints through a Use Case Diagram. The second is to test the checkout process more closely through an Activity Diagram. In both cases, social and ethical responsibilities are treated as part of the design rather than as afterthoughts. The Use Case Diagram defines system scope, actor relationships, mandatory dependencies, and optional extensions using standard UML conventions (OMG, 2017; Dennis, et al., 2015). The Activity Diagram then develops the checkout flow further by showing shipping choice, stock locking, fraud review, payment authorisation, and the eventual success or failure outcome (Fowler, 2004).

\begin{figure}[htbp]
\centering
\includegraphics[width=0.88\textwidth,trim=0cm 5.2cm 0cm 0.0cm,clip,page=1]{UML.drawio.pdf}
\caption{Amazon Enterprise Online POS Use Case Diagram}
\end{figure}
\section{The OO Analysis and Design with UML Techniques}

\subsection{The Design of Functional Requirements with Use Case Diagrams}

Figure~1 represents a platform that goes well beyond a basic shopping-basket application. On the customer side, it covers product search, cart handling, personalised recommendations, checkout, returns, privacy controls, accessibility settings, and optional features such as gift wrapping, promo gift cards, MFA, and carbon-neutral shipping. It also brings seller and administrative functions into the same model, since listing quality, inventory consistency, fraud management, and reporting all influence how the enterprise POS operates. This broader scope follows an important OOAD idea: the system boundary should reflect every outside entity that can trigger behaviour or shape outcomes, not just the primary purchaser (Dennis, et al., 2015).

\begin{table}[t]
\centering
\footnotesize
\setlength{\tabcolsep}{3pt}
\renewcommand{\arraystretch}{1.04}
\caption{Requirement Analysis Table for Amazon Enterprise Online POS}
\begin{tabular}{|>{\raggedright\arraybackslash}p{2.15cm}|>{\raggedright\arraybackslash}p{3.35cm}|>{\raggedright\arraybackslash}p{4.15cm}|>{\raggedright\arraybackslash}p{2.1cm}|>{\raggedright\arraybackslash}p{2.35cm}|}
\hline
\rowcolor[gray]{0.92}\textbf{List of Actors (Including Human Users, Machine Agents)} & \textbf{Functional Requirements} & \textbf{Rules/Conditions} & \textbf{Data Elements} & \textbf{Critique Notes} \\ \hline
Guest Customer (Human) & \textit{Search \& Filter Products}; \textit{Manage Shopping Cart} & Guest access covers browsing and cart use only; no direct link to checkout is shown. & Search query, product details, cart session & Fits Figure~1 and keeps guest scope credible. \\ \hline
Registered Customer (Human) & \textit{View Personalized Recs}; \textit{Proceed to Checkout}; \textit{Initiate Return Request}; \textit{Manage Data Privacy GDPR}; \textit{Set Accessibility Prefs} & Checkout includes shipping verification, fraud detection, and secure payment; privacy and accessibility remain visible user services. & Customer profile, basket data, address data, privacy settings, recommendation data & Strong actor-use-case alignment. \\ \hline
Third-Party Seller (Human) & \textit{Manage Vendor Listings}; \textit{Sync Vendor Inventory} & Seller stock and listing data must remain consistent with platform inventory. & Listing record, SKU, stock quantity & Clear separation from customer actions. \\ \hline
System Admin (Human) & \textit{Perform Fraud Detection}; \textit{Update Product Catalog}; \textit{Generate Sales Reports} & Admin links represent oversight, platform maintenance, and reporting rather than direct shopping activity. & Fraud indicators, catalogue data, sales metrics & Matches the actual diagram links. \\ \hline
Inventory System (Machine) & Supports \textit{Proceed to Checkout}; \textit{Sync Vendor Inventory}; \textit{Process Refund} & Stock must be checked, synchronised, and updated when refunds change order state. & SKU, stock level, reservation state & Correctly tied to checkout, refund, and vendor sync. \\ \hline
Payment Gateway (Machine) & Supports \textit{Process Secure Payment} & Payment authorisation is external and uses protected data exchange. & Payment token, transaction result, authorisation code & Keeps financial risk outside the POS core. \\ \hline
Logistics API (Machine) & Supports \textit{Verify Shipping Address} and optional green shipping & Shipping validation depends on logistics data; carbon-neutral shipping is optional. & Address record, delivery option, carbon-shipping flag & Green shipping extends address verification in Figure~1. \\ \hline
Analytics Engine (Machine) & Supports \textit{Generate Sales Reports} & Reporting is analytical rather than part of the live transaction path. & Report input, sales metrics, generated summary & Useful support system, but not part of checkout. \\ \hline
\end{tabular}
\end{table}

Checkout remains the central use case in the model. It \texttt{<<include>>}s \textit{Verify Shipping Address}, \textit{Perform Fraud Detection}, and \textit{Process Secure Payment}, because those steps must be completed before a purchase can finish successfully (OMG, 2017). Optional behaviour is shown through \texttt{<<extend>>}. In this design, \textit{Request Gift Wrapping} extends checkout, \textit{Apply Promo Gift Card} extends secure payment, \textit{Select Carbon-Neutral Shipping} extends shipping verification, and \textit{Enable Multi-Factor Authentication (MFA)} extends privacy management. \textit{Process Refund} is linked to both secure payment and vendor inventory synchronisation, which makes sense because a refund affects both financial records and stock status.

\smallskip\noindent\textbf{Use case specification: Proceed to Checkout.}\quad\textit{Primary actor:} Registered Customer.\;\textit{Supporting actors:} Inventory System, Payment Gateway, Logistics API.\;\textit{Preconditions:} customer authenticated; basket non-empty.\;\textit{Main flow:} (1)~confirm basket and select shipping; (2)~verify shipping address; (3)~perform fraud detection; (4)~process secure payment; (5)~confirm order and send notification.\;\textit{Alternative flows:} gift wrapping requested \textrightarrow\ add before payment; high fraud risk \textrightarrow\ escalate to manual review; payment declined \textrightarrow\ release stock and notify customer.\;\textit{Postconditions:} order confirmed with stock reserved, or transaction cancelled with stock released.
\smallskip

A number of constraints are made explicit in the model. Checkout is limited to registered customers, address verification depends on the Logistics API, inventory consistency is maintained through the Inventory System, and payment processing is handled through the Payment Gateway using tokenised channels aligned with PCI DSS v4.0 (PCI SSC, 2022). Accessibility preferences are included as a necessary design feature in line with WCAG~2.1 guidance (W3C, 2018). Fraud handling also reflects the view that high-stakes automated decisions should remain open to human review (ACM, 2018). Presenting privacy, MFA, and related governance measures as visible use cases rather than hidden background assumptions helps the model respond directly to the brief's requirement to identify assumptions, conditions, constraints, and social or ethical considerations.


Taken as a whole, Figure~1 is both complete and defensible. Every actor is connected to a recognisable service, the \texttt{<<include>>} and \texttt{<<extend>>} relationships are used for clear reasons, and the external systems are tied to genuine dependencies rather than decorative additions (Booch, et al., 1999). The main limitation is that the scope is relatively wide: buying, returns, compliance, seller operations, and administrative reporting all appear within one diagram. In this case, that breadth is still manageable because the use cases remain readable and cluster naturally by role. Using standard UML validation criteria, the model can be regarded as \textit{complete} (all identified actors have at least one goal), \textit{accurate} (the relationships reflect real system dependencies), and \textit{aligned with user requirements} (the functional scope, constraints, and ethical concerns can be traced from Table~1 to Figure~1).

\begin{figure}[htbp]
\centering
\includegraphics[width=0.72\textwidth,trim=0cm 9.6cm 0cm 0.6cm,clip,page=1]{ActivityDiagram.drawio.pdf}
\caption{Activity Diagram for Checkout Confirmation, Shipping, and Payment}
\end{figure}
\subsection{The Design of Functional Behaviour with Activity Diagrams}

After reviewing the use case model in Section~2.1, \textit{Proceed to Checkout} was chosen for activity modelling because it represents the most critical transaction path in the system. It combines customer actions, inventory management, payment authorisation, and ethical control points within one workflow (Sommerville, 2016). Figure~2 develops that process in more detail. The customer first confirms the order and selects a shipping option. When green shipping is chosen, the Inventory and Logistics lane calculates carbon credits before the flow returns to stock verification and locking. If standard shipping is selected instead, that additional step is bypassed.


Within the Amazon POS lane, the system checks the order details and then splits into two parallel controls: an AI-based fraud risk scan and PCI DSS payment-data encryption (PCI SSC, 2022). These branches join again before the decision point labelled \textit{Fraud Risk?}. Where the risk is low, the transaction moves directly to payment authorisation. Where the risk is high, the flow is redirected to \textit{Human Oversight: Manual Review} and then to \textit{Approve?}, reflecting the principle that financial decisions of this kind should not be left entirely to automation (ACM, 2018). The model also introduces a 10-minute timeout around the bank-handshake stage so that the transaction cannot remain pending without limit.

The ending behaviour is designed to be unambiguous. After \textit{Request Payment Auth}, the flow checks \textit{Payment OK?}. If the payment succeeds, the system generates an invoice and sends confirmation to the customer. If the payment fails, or if the transaction is rejected during manual review, the system releases the reserved stock, cancels the order, and issues a failure alert. This clear termination logic ensures that every path leads to a definite final state, which is an important requirement for completeness in UML activity modelling (Fowler, 2004). Figure~2 therefore develops Figure~1 in a consistent way: checkout is broken down into shipping choice, fraud control, secure payment, and a clear success/failure split, while ethical oversight remains visible within the flow. On that basis, UML validation suggests that the activities, actors, decision logic, guard conditions, and business rules identified earlier in the use case analysis are represented correctly and fully in the activity model (OMG, 2017).

\section{Conclusions and Recommendations}

This report has used OOAD and UML to define the scope of an Amazon enterprise online POS system and to examine its most important transaction path (Booch, et al., 1999). The use case model sets out the responsibilities of customers, sellers, administrators, and machine agents in a structured and readable form (Dennis, et al., 2015). The activity model then shows how shipping choice, stock control, fraud review, payment authorisation, and customer notification operate together during checkout. Overall, the design indicates that an enterprise POS should balance usability with security, accessibility, and ethical oversight, in line with current standards and professional guidance (PCI SSC, 2022; W3C, 2018; ACM, 2018).

\textbf{Recommendations for future improvement:}
\begin{enumerate}
\item \textbf{Add explicit return-case behaviour modelling:} the use case diagram already includes \textit{Initiate Return Request} and \textit{Process Refund}, so a later stage should add an activity model for returns, refund approval, and stock reinstatement.
\item \textbf{Strengthen admin-side fraud reporting:} the current models imply oversight, but a future revision could add a clearer reporting or alert workflow that shows how suspicious transactions are surfaced to administrators and audited over time.
\end{enumerate}

\begin{thebibliography}{9}
\small
\bibitem{acm2018} ACM (2018) \textit{ACM Code of Ethics and Professional Conduct}. Available at: \url{https://www.acm.org/code-of-ethics} (Accessed: 25 March 2026).
\bibitem{booch1999} Booch G., Rumbaugh J. and Jacobson I. (1999) \textit{The Unified Modeling Language User Guide}, Addison-Wesley, ISBN: 0-201-57168-4.
\bibitem{dennis2015} Dennis A., Wixom B.H. and Tegarden D. (2015) \textit{Systems Analysis and Design with UML}, 5th edn., John Wiley \& Sons, ISBN: 978-1-118-80436-1.
\bibitem{fowler2004} Fowler M. (2004) \textit{UML Distilled: A Brief Guide to the Standard Object Modeling Language}, 3rd edn., Addison-Wesley, ISBN: 0-321-19368-7.
\bibitem{omg2017} Object Management Group (2017) \textit{OMG Unified Modeling Language (UML) Version 2.5.1}. Available at: \url{https://www.omg.org/spec/UML/2.5.1} (Accessed: 25 March 2026).
\bibitem{pcissc2022} PCI Security Standards Council (2022) \textit{Payment Card Industry Data Security Standard (PCI DSS) v4.0}. Available at: \url{https://www.pcisecuritystandards.org} (Accessed: 25 March 2026).
\bibitem{sommerville2016} Sommerville I. (2016) \textit{Software Engineering}, 10th edn., Pearson, ISBN: 978-1-292-09613-1.
\bibitem{w3c2018} W3C (2018) \textit{Web Content Accessibility Guidelines (WCAG) 2.1}. Available at: \url{https://www.w3.org/TR/WCAG21/} (Accessed: 25 March 2026).
\end{thebibliography}

\clearpage
\section*{Effort Allocation Sheet}

\begin{table}[htbp]
\centering
\small
\begin{tabularx}{\textwidth}{|p{2.2cm}|X|c|c|p{1.8cm}|}
\hline
\rowcolor[gray]{0.92}\textbf{Name} & \textbf{Individual Contributions} & \textbf{Effort (\%)} & \textbf{Reflection (Y/N)} & \textbf{Signature} \\ \hline
Fu Renjie (32804061) & Introduction, requirement analysis, coordination, report integration & 100\% & Y & \includegraphics[angle=90,width=1.2cm]{70d35bbee1d874cb9afc1c6688761d50.jpg} \\ \hline
Wu Renxuan (32803720) & Use Case Diagram, use case validation, static-model consistency checking & 100\% & Y & \includegraphics[angle=90,width=1.2cm]{59da8176fa6de3c3e3eca260ecbd4da8.jpg} \\ \hline
Cui Zhiqing (32803755) & Activity Diagram, behavioural validation, conclusion and recommendations & 100\% & Y & \includegraphics[width=1.5cm]{10ef88f5bd748f2e19d2fbe40340d64b.jpg} \\ \hline
\end{tabularx}
\end{table}

\section*{Appendix: Individual Reflections}

\textbf{Reflection by Fu Renjie (Coordination \& Analysis):} \\
This project made me realise that the most difficult part often appears right at the beginning. When we first started discussing the Amazon POS domain, nearly every feature seemed important enough to keep. We talked about payment, logistics, privacy, recommendations, seller functions, and accessibility, and it was not easy to decide how much of that belonged in Stage~1. My role was to coordinate the group and support the requirement analysis, so I spent a lot of time helping us decide what should remain inside the scope and what should be left for later work.

What I learned most clearly was the importance of consistency across the whole report. Once one part of the model changed, the table, the written discussion, and the second diagram often needed to be updated as well. That experience changed the way I think about coordination. It is not only about sharing tasks or making sure deadlines are met. It is also about keeping the overall submission coherent. Our draft reviews and group feedback helped a lot here because they exposed places where our explanation of constraints and ethical concerns was still too weak.

I used AI tools near the end of the process to improve sentence flow and to check whether key terms were being used consistently across sections. The actual decisions about scope and model logic still came from our group discussions. Overall, this coursework improved my confidence in planning, abstraction, and coordinating work within a team.

\bigskip
\textbf{Reflection by Wu Renxuan (Static Modelling):} \\
My main responsibility was the Use Case Diagram and the requirement analysis table. At the start, I found it surprisingly difficult to distinguish between mandatory behaviour and optional behaviour. I initially treated \texttt{<<include>>} and \texttt{<<extend>>} as if they were mainly diagram layout choices. After revising the diagram several times and talking through it with my teammates, I understood the distinction much better. Include shows that the main use case depends on that action, while extend is used when extra behaviour appears only in a particular situation.

That changed how I approached the checkout model. I became more careful about deciding which relationships were essential and which should remain optional. I also learned that a use case diagram cannot really stand alone. The actor names, use case labels, and written explanation all need to support one another. If they do not line up, the report feels less convincing even when the diagram itself is technically acceptable.

Earlier in the process I used AI to look at a few examples of UML use case relationships. That helped me frame better questions, but it did not make the final decisions for us. The final version of the model came from repeated revision and group discussion. Through this part of the coursework, I developed a stronger understanding of UML accuracy, naming, and consistency between model and explanation.

\bigskip
\textbf{Reflection by Cui Zhiqing (Dynamic Modelling):} \\
The activity diagram was the part of the report that I revised most often. At first I assumed the checkout flow would be straightforward to present, but it quickly became clear that the process involved more detail than I expected. We needed to show order validation, stock locking, fraud screening, payment authorisation, timeout handling, and two different ending states. Working through that made me understand that an activity diagram is much more than a simple flowchart. It needs a clear sequence, clear decisions, and clear responsibility for each step.

Another lesson for me was the importance of layout. In the early drafts, the diagram felt crowded and some paths were difficult to follow. After experimenting with different swimlane arrangements, I found that separating the POS core logic, the Inventory and Logistics lane, and the End Customer lane made the whole process easier to read. I used AI at this stage to test wording and layout ideas, but the final flow and design choices still came from our own discussions.

The strongest thing I learned was the need to check the activity model against the use case model carefully. If checkout included a behaviour in the use case analysis, the activity diagram also needed to represent it clearly. Our review sessions helped us catch missing guard conditions and one incorrect flow direction. Correcting those issues improved the final model and showed me how important careful checking is in behavioural design.

\end{document}
